
% INTRODUÇÃO
\chapter{INTRODUÇÃO}
\label{chap:introducao}

O presente relatório descreve as atividades desenvolvidas durante o estágio supervisionado realizado entre fevereiro e junho de 2026 no Laboratório ASSERT, vinculado ao Polo de Inovação do Instituto Federal da Paraíba (IFPB). As atividades concentraram-se na evolução do sistema Banco de Especialistas, plataforma destinada ao gerenciamento de especialistas e ao apoio à composição de equipes para projetos de Pesquisa, Desenvolvimento e Inovação (PD&I).

O Banco de Especialistas permite o cadastro, a classificação e a avaliação de candidatos com base em critérios como formação acadêmica, experiência profissional e participação em projetos. Nesse contexto, a plataforma reúne informações utilizadas na gestão dos especialistas e nos processos de seleção e composição de equipes vinculadas às iniciativas do Polo de Inovação.

Durante o estágio, identificou-se uma demanda relacionada à alteração de categoria dos especialistas. O procedimento então utilizado exigia que o interessado solicitasse a mudança por e-mail e que um administrador realizasse manualmente uma consulta SQL diretamente no banco de dados. A ausência de um fluxo formal de solicitações dificultava o acompanhamento das alterações e a recuperação do estado anterior dos registros, além de aumentar o risco de inconsistências decorrentes de intervenções diretas na base de dados.

Diante desse cenário, o trabalho concentrou-se no desenvolvimento de um módulo de solicitações para formalizar e gerenciar operações de alteração cadastral e de classificação dos especialistas. A solução contempla mecanismos para registrar solicitações, acompanhar seu andamento, permitir sua análise e decisão e manter informações necessárias à rastreabilidade das operações. Dessa forma, busca-se reduzir a dependência de intervenções manuais no banco de dados e estabelecer um processo mais controlado e verificável.

Do ponto de vista técnico, o desenvolvimento envolveu conceitos de engenharia de software, aplicações web, persistência de dados relacionais, testes automatizados e conteinerização. Destacam-se a integração entre componentes de front-end e back-end, o uso de mecanismos de registro e recuperação do estado dos usuários e a manutenção de um histórico das operações. Esses fundamentos são apresentados no capítulo dedicado ao embasamento teórico.

\section{Objetivos}
\label{sec:objetivos}

\subsection{Objetivo geral}
\label{subsec:objetivo-geral}

Desenvolver um módulo de solicitações de alteração cadastral e de classificação no sistema Banco de Especialistas, contemplando os fluxos de criação, consulta, aprovação, rejeição e reversão, com mecanismos de rastreabilidade e registro das operações.

\subsection{Objetivos específicos}
\label{subsec:objetivos-especificos}

\begin{itemize}
    \item Levantar e analisar os requisitos e as regras de negócio relacionados às solicitações de alteração cadastral e de classificação dos especialistas.
    
    \item Modelar os componentes de persistência necessários ao registro das solicitações e à recuperação do estado dos usuários.
    
    \item Implementar as funcionalidades de gerenciamento das solicitações no back-end, incluindo as regras de negócio, a integração com o PostgreSQL e o acesso aos dados por meio do SQLAlchemy.
    
    \item Desenvolver as interfaces web necessárias à criação, ao acompanhamento e à análise das solicitações.
    
    \item Implementar mecanismos de rastreabilidade das operações, incluindo o registro de eventos e a apresentação de seu histórico.
    
    \item Desenvolver e executar testes unitários e de integração para verificar as regras de negócio e o funcionamento integrado dos componentes da aplicação.
\end{itemize}

\section{Escopo e limitações}
\label{sec:escopo-limitacoes}

O trabalho abrange o desenvolvimento do módulo de solicitações e as funcionalidades necessárias ao gerenciamento das operações de alteração cadastral e de classificação dos especialistas no sistema Banco de Especialistas. Estão incluídos os mecanismos de registro, acompanhamento, decisão e reversão das solicitações, bem como os recursos de persistência e rastreabilidade associados a esse fluxo.

Não fazem parte do escopo a redefinição dos critérios de pontuação dos especialistas, a modificação do processo seletivo, a substituição do sistema de avaliação existente e a implementação de notificações por e-mail. Também não é contemplado o tratamento de múltiplas solicitações simultâneas para um mesmo especialista. Essas limitações delimitam as responsabilidades da solução desenvolvida e podem orientar evoluções futuras do sistema.

\section{Procedimentos metodológicos}
\label{sec:procedimentos-metodologicos}

O trabalho caracteriza-se como um relato técnico de natureza aplicada e caráter descritivo, desenvolvido no contexto de um estágio supervisionado em desenvolvimento de software. Sua execução compreendeu o estudo da arquitetura e das regras de negócio do sistema, o levantamento e a análise dos requisitos, a modelagem dos componentes necessários, a implementação das funcionalidades e a realização de testes.

A verificação técnica envolveu testes unitários e de integração, destinados a avaliar o comportamento das funcionalidades implementadas e a interação entre os componentes da aplicação. Também foi verificado o funcionamento integrado do front-end e do back-end em ambiente de desenvolvimento conteinerizado com Docker. As atividades foram conduzidas no contexto da arquitetura e das tecnologias já adotadas pelo projeto, considerando suas restrições e regras de negócio.

\section{Contexto institucional}
\label{sec:contexto-institucional}

O estágio foi realizado no âmbito do Polo de Inovação do Instituto Federal da Paraíba, estrutura associada à Agência de Inovação da instituição (NEO-IFPB). O Polo atua no desenvolvimento de iniciativas de Pesquisa, Desenvolvimento e Inovação (PD&I), aproximando a instituição de demandas tecnológicas do setor produtivo e da sociedade.

As atividades técnicas foram desenvolvidas no Laboratório ASSERT, ambiente vinculado ao Polo de Inovação no qual são realizadas atividades de desenvolvimento e manutenção de sistemas computacionais para projetos institucionais. Nesse contexto, o estágio proporcionou a participação nas etapas de análise, implementação, integração e verificação de funcionalidades de uma aplicação web existente.

\section{Organização do relatório}
\label{sec:organizacao-relatorio}

Além deste capítulo introdutório, o relatório está organizado em três capítulos. O segundo apresenta o embasamento teórico relacionado às tecnologias e aos conceitos utilizados no desenvolvimento. O terceiro descreve as atividades realizadas durante o estágio, incluindo a análise dos requisitos, a implementação das funcionalidades e os testes executados. Por fim, o quarto capítulo apresenta as considerações finais, sintetizando os resultados do trabalho, as limitações identificadas e os possíveis desdobramentos.
